\begin{figure}[H]
  \centering
  \resizebox{0.68\textwidth}{!}{%
  \begin{tikzpicture}[node distance=6mm and 9mm]
    \node[casePrimary,minimum width=52mm] (tauri)
      {\texttt{build desktop} / \texttt{tauri build}};
    \node[caseBox,minimum width=52mm,below=of tauri] (strategy)
      {eine Zielstrategie wählen};
    \node[caseBox,minimum width=32mm,below left=of strategy] (windows) {Windows-Host};
    \node[caseBox,minimum width=32mm,below=of strategy] (macos) {macOS-Host};
    \node[caseBox,minimum width=32mm,below right=of strategy] (linux) {Linux-Host};
    \node[caseBox,minimum width=58mm,below=15mm of macos] (packaging)
      {Tauri-Packaging\\native Paketformate je Plattform};
    \node[caseBox,minimum width=58mm,below=of packaging] (artifacts)
      {unsignierte Bundle-Artefakte};
    \node[caseOptional,minimum width=82mm,below=of artifacts] (external)
      {Signing / macOS-Notarisierung / Veröffentlichung / Updater\\geschützte produktspezifische Integration};

    \draw[caseFlow] (tauri) -- (strategy);
    \draw[caseFlow] (strategy) -- (windows);
    \draw[caseFlow] (strategy) -- (macos);
    \draw[caseFlow] (strategy) -- (linux);
    \draw[caseFlow] (windows) |- (packaging);
    \draw[caseFlow] (macos) -- (packaging);
    \draw[caseFlow] (linux) |- (packaging);
    \draw[caseFlow] (packaging) -- (artifacts);
    \draw[caseSecondaryFlow] (artifacts) -- (external);
  \end{tikzpicture}}
  \caption{Plattformübergreifendes Modell für Desktop-Releases}
  \label{fig:desktop-release-model}
  \smallskip
  {\footnotesize Quelle: Eigene Darstellung auf Grundlage von
  \textcite{kleiveist2026release}.\par}
\end{figure}
